\section*{第 \ref{ch:g12:doppler-effect} 章——多普勒效应}

\begin{solution}{exo:g12:doppler-effect:1}
(a) $f' = 700 \times 340/315 \approx \qty{756}{Hz}$。
(b) $f' = 700 \times 340/365 \approx \qty{652}{Hz}$。
\qty{700}{Hz}：车与警笛一起运动，二者间距始终不变。
\end{solution}

\begin{solution}{exo:g12:doppler-effect:2}
\emph{响度}减弱是因为波向外扩散——它取决于距离。
\emph{音调}由接近转为远离时降低——这是多普勒效应，它取决于
距离变化的快慢。两种效应相互独立。
\end{solution}

\begin{solution}{exo:g12:doppler-effect:3}
运动的观察者：$f' = 700\,(1 + 6.0/340) \approx \qty{712}{Hz}$；远离警笛骑行时，
$700\,(1 - 6.0/340) \approx \qty{688}{Hz}$。警笛相对空气静止，因此波
形不受扰动；只有观察者穿过它。
\end{solution}

\begin{solution}{exo:g12:doppler-effect:4}
前方：$\lambda' = 310/400 = \qty{0.775}{m}$；后方：$370/400 = \qty{0.925}{m}$；
静止时：$340/400 = \qty{0.85}{m}$。短波长传给列车所驶向的人，
较长的则传给列车所驶离的人。
\end{solution}

\begin{solution}{exo:g12:doppler-effect:5}
$\Delta f/f \approx 5.0/340 \approx 1.5\%$：$\Delta f \approx \qty{6.5}{Hz}$，因此约
\qty{446}{Hz}——相当于四分之一个半音。音乐家刚好能察觉。
\end{solution}

\begin{solution}{exo:g12:doppler-effect:6}
墙面接收到的频率为 $f\,v/(v - v_b)$（运动声源）；蝙蝠作为运动的观察者，
正接近回声，听到的频率是其 $(1 + v_b/v)$ 倍：$f' = f(v + v_b)/(v - v_b) = 50.0 \times
346/334 \approx \qty{51.8}{kHz}$，频移为 $+\qty{1.8}{kHz}$。这是雷达测速枪
的双重频移，只是发射源与反射体的角色互换：只有接近速率起作用。
\end{solution}

\begin{solution}{exo:g12:doppler-effect:7}
$u = c\,\Delta f/(2f) = \num{3.00e8} \times 4000/(2 \times \num{24.15e9}) \approx
\qty{24.8}{m/s} \approx \qty{89}{km/h}$。当速率为 \qty{50}{km/h}（$\qty{13.9}{m/s}$）时：
$\Delta f = 2 \times 13.9 \times \num{24.15e9}/\num{3.00e8} \approx \qty{2.2}{kHz}$。
\end{solution}

\begin{solution}{exo:g12:doppler-effect:8}
$u = v\,\Delta f/(2f) = 1540 \times 2600/\num{1.0e7} \approx \qty{0.40}{m/s}$。
与超声束垂直的流动，径向速度为零，因此完全没有频移。
\end{solution}

\begin{solution}{exo:g12:doppler-effect:9}
(a) 声源：$f' = 1000 \times 340/306 \approx \qty{1111}{Hz}$。(b) 观察者：
$f' = 1000 \times 1.1 = \qty{1100}{Hz}$。慢速运动预言为 \qty{1100}{Hz}，对观察者精确，
对声源则偏小 \qty{11}{Hz}——即 $x^2$ 阶修正，为 $1\%$
的 $f$（$x = 0.1$ 时）。
\end{solution}

\begin{solution}{exo:g12:doppler-effect:10}
$f = 2 \times 465 \times 415/880 \approx \qty{439}{Hz}$；
$v_s = 340 \times 50/880 \approx \qty{19.3}{m/s} \approx \qty{70}{km/h}$。真实的 $f$
是调和平均：接近时频率升高的幅度大于远离时降低的幅度，因此
$f$ 低于算术平均值（\qty{440}{Hz}）。
\end{solution}

\begin{solution}{exo:g12:doppler-effect:11}
波长变长：红移，恒星远离。$v_r = c\,\Delta\lambda/\lambda =
\num{3.00e8} \times 0.2/656.3 \approx \qty{9.1e4}{m/s} \approx \qty{91}{km/s}$。
\end{solution}

\begin{solution}{exo:g12:doppler-effect:12}
$\Delta\lambda = \lambda K/c = 500 \times 55/\num{3.00e8} \approx \qty{9.2e-5}{nm}$；
$\Delta\lambda/\lambda = 55/\num{3.00e8} \approx \num{1.8e-7}$。整个信号仅为
千万分之二：光谱仪漂移超过 $10^{-7}$ 就会将其淹没。
\end{solution}

\begin{solution}{exo:g12:doppler-effect:13}
$v_r = 0.024\,c \approx \qty{7.2e6}{m/s}$。速率与距离成正比，
各个方向都相同：所有距离均匀伸展——这是一个没有特殊中心的膨胀宇宙。
在 $2.4\%$ 时，修正量的数量级为 $(v_r/c)^2 \approx
\num{6e-4}$：仍然可靠。
\end{solution}

\begin{solution}{exo:g12:doppler-effect:14}
两侧边缘向相反方向偏移，因此圆盘平均后的谱线是\emph{展宽}的，而不是
整体偏移：全宽为 $2\lambda v/c = 2 \times 656.3 \times \num{2.0e3}/\num{3.00e8}
\approx \qty{8.8e-3}{nm}$（约 $\pm\qty{4.4e-3}{nm}$）。
\end{solution}

\begin{solution}{exo:g12:doppler-effect:15}
(a) 车前方的波峰间距为 $(v - v_s)T$；卡车以速率 $v_o$ 远离，
并以相对速率 $v - v_o$ 与其相遇：$f' = (v - v_o)/\bigl((v - v_s)T\bigr) =
f(v - v_o)/(v - v_s)$。(b) $f' = 700 \times 310/300 \approx \qty{723}{Hz}$。
(c) 当 $v_o = v_s$ 时 $f' = f$：间距保持恒定，
距离不变，按定义即无多普勒频移。
\end{solution}

\begin{solution}{pb:g12:doppler-effect:1}
\textbf{1.} 每个波峰都从发出它的位置向外扩展：波峰在移动的喇叭前方挤在一起，
在其后方被拉开，因此接近的手机接收到的
波峰频率更高（$f_1$），高于远离的手机（$f_2$）。

\textbf{2.} $f_1 = f\,\dfrac{v}{v - u}$，\quad $f_2 = f\,\dfrac{v}{v + u}$。

\textbf{3.} $f_1 - f_2 = \frac{2fv\,u}{v^2 - u^2}$ 和 $f_1 + f_2 =
\frac{2fv^2}{v^2 - u^2}$；相除得 $u = v\,\frac{f_1 - f_2}{f_1 + f_2} =
340 \times 60/882 \approx \qty{23}{m/s} \approx \qty{83}{km/h}$。

\textbf{4.} $\dfrac{2f_1 f_2}{f_1 + f_2} = f$（公因子约去）：
$f = 2 \times 471 \times 411/882 \approx \qty{439}{Hz}$——调和平均，低于
算术平均值 \qty{441}{Hz}，因为接近时使 $f$ 升高的幅度大于远离时使其降低的幅度。

\textbf{5.} 前方：$(v - u)/f = 316.9/439 \approx \qty{0.72}{m}$；后方：
$363.1/439 \approx \qty{0.83}{m}$；静止时：$340/439 \approx \qty{0.77}{m}$。

\textbf{6.} 一次作为接收波的运动\emph{观察者}，一次作为
转发所接收波的运动\emph{波源}。

\textbf{7.} $f' = f\,(1 + u/c)/(1 - u/c) \approx f(1 + 2u/c)$，因此
$\Delta f = f' - f \approx \dfrac{2u}{c}\,f$。

\textbf{8.} $u = c\,\Delta f/(2f) = \num{3.00e8} \times 3550/(2 \times \num{24.125e9})
\approx \qty{22}{m/s}$。

\textbf{9.} $22.1 \times 3.6 \approx \qty{79}{km/h}$：刚好低于 \qty{80}{km/h}
限速。

\textbf{10.} $\Delta f = 2f/(3.6\,c) \approx \qty{45}{Hz}$ 每 \unit{km/h}。将回声与发出的波
相拍，只剩下\emph{差频}：一个音频频率，
计数十分容易，尽管它只是载波的 $10^{-7}$ 分之一。

\textbf{11.} 夜复一夜测量已知谱线的波长；随后
由 $v_r = c\,\Delta\lambda/\lambda$ 求得。

\textbf{12.} $K = \qty{55}{m/s}$，$P = \qty{4.2}{days}$。

\textbf{13.} $\Delta\lambda = 550 \times 55/\num{3.00e8} \approx \qty{1.0e-4}{nm}$，
$\Delta\lambda/\lambda \approx \num{1.8e-7}$——即五百万分之一，约小
$\num{4e5}$ 倍，与火车的 $7\%$ 频移相比。

\textbf{14.} 恒星与行星相互吸引、大小相等，因此二者都绕共同的
质心运动：恒星不可能静止不动。正弦形的 $v_r$ 是匀速圆周运动，
从侧面看去——一条圆轨道。

\textbf{15.} 只有视线方向的分量会改变恒星与地球之间的距离，从而
改变波长。若正对观察，则 $v_r = 0$ 始终成立：完全没有信号。

\textbf{16.} $P = \qty{3.63e5}{s}$：$r^3 = \num{6.67e-11} \times \num{2.0e30} \times
(\num{3.63e5})^2/(4\pi^2) \approx \num{4.4e29}$，因此 $r \approx \qty{7.6e9}{m}$——
约为日地距离的 $1/20$。

\textbf{17.} $V = 2\pi r/P = 2\pi \times \num{7.6e9}/\num{3.63e5} \approx
\qty{1.3e5}{m/s}$。

\textbf{18.} $m = M K/V = \num{2.0e30} \times 55/\num{1.3e5} \approx \qty{8.3e26}{kg}$。

\textbf{19.} 质量约为木星的 $0.4$ 倍（$\approx 140$ 个地球质量），
轨道半径比水星小 $8$ 倍：一颗紧贴恒星炙烤的气态巨行星——“热木星”。

\textbf{20.} 轨道倾角遮掉了一部分运动：所测得的 $K$ 只是径向分量，
因此真实质量\emph{更大}——\qty{8.3e26}{kg} 为最小值。一次周期性的
仅为五百万分之一的频移，就给出了一颗
从未有人见过的行星的轨道、速率与最小质量。
\end{solution}
